\centerline{\bf \Huge Хата скраю в графах}
\centerline{\bf \Huge }

\begin{flushright}
        \scriptsize{}
\end{flushright}

\problem{1}
Баир нарисовал карту из нескольких городов и дорог между ними (дорог получилось не менее двух, каждая дорога соединяет различные города). На каждом городе он написал, сколько дорог из него выходит. Могло ли получиться так, что на этой карте нет двух дорог с совпадающими наборами чисел на концах?

\problem{2}
В компании из ста тысяч человек среди любых десяти есть трое попарно знакомых. Докажите, что можно выбрать восьмерых из них так, чтобы любой из оставшихся был знаком с кем-то из этих восьмерых.

\problem{3}
В деревне дураков функционирует несколько клубов анонимных алкоголиков. Каждый житель деревни входит хотя бы в $k$ клубов. Любые два клуба содержат как максимум одного общего жителя. Докажите, что найдется не менее $k$ клубов с одинаковым числом участников.

\problem{4}
$\bigl[\textbf{Финал ВсОШ, 2004, 9.6}\bigr]$
В кабинете президента стоят 2004   телефона, любые два из которых соединены проводом одного из четырех цветов. Известно, что провода всех четырех цветов присутствуют. Всегда ли можно выбрать несколько телефонов так, чтобы среди соединяющих их проводов встречались провода ровно трех цветов?